\begin{textsample}{POS Dim 1 – ac – Score 44.00 – ac/ac\_em\_7.txt}  \label{ex:f1_pos_178}
6. \textbf{Competências} do \textbf{século} XXI Não basta que cada indivíduo acumule no começo da vida uma determinada quantidade de conhecimentos de que possa abastecer -se indefinidamente.
É, antes, necessário estar à altura de aproveitar e explorar, do começo ao fim da vida, todas as ocasiões de atualizar, aprofundar e enriquecer seus conhecimentos e de se adaptar a um mundo em mudança.
Nesse sentido, o relatório para a UNESCO, da Comissão Internacional de Educação para o \textbf{século} XXI ( 1998 ), é um documento que aponta discussões e orientações para a educação na busca do desenvolvimento dos países, a paz e a superação de \textbf{problemas} gerados num mundo que se desenvolve de \textbf{maneira} rápida e a proporções gigantescas.
Dentre as reflexões destacadas pelo relatório ( Delors, 1998 ), a educação deve se organizar em torno de quatro aprendizagens fundamentais que, ao longo de toda a vida, serão, de algum modo, para cada indivíduo, os pilares do conhecimento : - Aprender a conhecer, combinando uma cultura geral, suficientemente vasta, com a possibilidade de trabalhar em \textbf{profundidade} um pequeno número de matérias.
O que também significa : aprender a aprender, para beneficiar -se das oportunidades oferecidas pela educação ao longo de toda a vida. - Aprender a fazer, a fim de adquirir, não somente uma qualificação profissional, mas, de uma \textbf{maneira} mais ampla, \textbf{competências} que tornem a pessoa apta a enfrentar numerosas situações e a trabalhar em equipe.
Mas também aprender a fazer, no âmbito das diversas experiências sociais ou de trabalho que se oferecem aos jovens e adolescentes, quer espontaneamente, fruto do contexto local ou nacional, quer formalmente, graças ao desenvolvimento do ensino alternado com o trabalho. - Aprender a viver juntos desenvolvendo a \textbf{compreensão} do outro e a percepção das interdependências — realizar projetos comuns e preparar -se para gerir conflitos — no respeito pelos valores do pluralismo, da \textbf{compreensão} mútua e da paz. - Aprender a ser, para melhor desenvolver a sua personalidade e estar à altura de agir com cada vez maior \textbf{capacidade} de autonomia, de discernimento e de responsabilidade pessoal.
Para isso, não negligenciar na educação nenhuma das potencialidades de cada indivíduo : memória, \textbf{raciocínio}, sentido estético, \textbf{capacidades} físicas, aptidão para comunicar -se.
No \textbf{século} XXI estes fenômenos assumiram ainda mais amplitude, pois de acordo com Delors ( 1998 ), mais do que preparar os estudantes para uma dada sociedade, o desafio será, então, fornecer -lhes constantemente estímulos e referências intelectuais que lhes permitam compreender o mundo ao seu redor e lhes ensine a se comportar nele como \textbf{atores} responsáveis e justos.
Mais do que nunca, a educação parece ter como papel essencial conferir a todos os seres humanos a liberdade de pensamento, discernimento, sentimentos e imaginação de que necessitam para desenvolver os seus talentos e permanecerem, \textbf{tanto} quanto possível, donos do seu próprio destino.
A partir desses trabalhos e de outras pesquisas, como a Taxonomia dos Objetivos Educacionais de Bloom, de 1956, foi possível identificar e organizar as descobertas sobre o comportamento e o pensamento humano.
Assim, para descrever as \textbf{competências} do \textbf{século} XXI, a NATIONAL RESEARCH COUNCIL ( 2012 ) identificou três domínios de \textbf{competência} : - Domínio cognitivo, que envolve o pensamento e habilidades relacionadas, como \textbf{raciocínio}, resolução de \textbf{problemas} e memória. - Domínio intrapessoal, como o domínio afetivo de Bloom, que envolve emoções e sentimentos e inclui autorregulação - a \textbf{capacidade} de definir e \textbf{atingir} seus objetivos. - Domínio interpessoal, que é usado \textbf{tanto} para expressar informações para os outros quanto para interpretar as mensagens dos outros ( verbais e não verbais ) e responder adequadamente.
As distinções entre os três domínios, segundo Almlund et al. ( 2011 ) se refletem em como eles são delineados, estudados e medidos.
No domínio cognitivo, o conhecimento e as habilidades são normalmente medidos com testes de \textbf{capacidade} cognitiva geral ( também conhecida como teste do QI ) ou com testes mais específicos com \textbf{foco} em componentes escolares ou conteúdo relacionado ao trabalho.
A pesquisa sobre \textbf{competências} intrapessoais e interpessoais geralmente usa medidas de traços gerais de personalidade ou do temperamento infantil – tendências comportamentais gerais, como atenção ou timidez.
É inegável que os investimentos na educação pública contribuem para o bem comum, \textbf{fomentam} a estabilidade de uma família, bairros, comunidades, cidades e da nação.
Os estudantes de hoje podem enfrentar desafios futuros se sua escolaridade e atividades de aprendizagem os preparam para papéis de adultos como cidadãos, funcionários, gerentes, pais, voluntários, empresários e políticos.
Para \textbf{atingir} todo o seu \textbf{potencial} como adultos, os jovens precisam desenvolver uma gama de habilidades e conhecimentos que facilitem o domínio e a aplicação de saberes nas áreas das Linguagens, Matemática, Ciências Humanas e Ciências da Natureza, ao mesmo tempo, nos eixos estruturantes de Investigação Científica, Processos Criativos, Mediação e Intervenção Sociocultural e Empreendedorismo.
Isso porque o contexto \textbf{atual} \textbf{exige} que se desenvolvam habilidades para a resolução de \textbf{problemas}, pensamento \textbf{crítico}, comunicação, colaboração e autogestão – mais conhecidas como"\textbf{Competências} do Século XXI"( NRC, 2012 ).
Para \textbf{tanto}, é necessário que os sistemas, redes e escolas garantam um patamar comum de aprendizagens a todos os estudantes, \textbf{tarefa} para a qual a BNCC é instrumento fundamental, uma vez que, ao longo da Educação Básica, as aprendizagens essenciais definidas devem concorrer para assegurar aos estudantes o desenvolvimento das dez \textbf{competências} gerais, que consubstanciam, no âmbito pedagógico, os direitos de aprendizagem e desenvolvimento.
A BNCC, mostra -se também alinhada ao objetivo 4 “ Assegurar a educação inclusiva e equitativa e de qualidade, e promover oportunidades de aprendizagem ao longo da vida para todos ” da \textbf{Agenda} 2030 da Organização das Nações Unidas ( BRASIL, 2018 ).
É \textbf{imprescindível} destacar que as \textbf{competências} gerais da Educação Básica, apresentadas a seguir, \textbf{inter-relacionam} -se e desdobram -se no tratamento didático proposto para as três etapas da Educação Básica ( Educação Infantil, Ensino Fundamental e Ensino Médio ), articulando -se na construção de conhecimentos, no desenvolvimento de habilidades e na formação de atitudes e valores, nos termos da LDB. \textbf{QUADRO} 1 | \textbf{COMPETÊNCIAS} GERAIS DA EDUCAÇÃO BÁSICA 1.
Conhecimento - Valorizar e utilizar os conhecimentos historicamente construídos sobre o mundo físico, social, cultural e \textbf{digital} para entender e explicar a realidade, continuar aprendendo e colaborar para a construção de uma sociedade justa, democrática e inclusiva. 2.
Pensamento científico, \textbf{crítico} e criativo - Exercitar a curiosidade intelectual e recorrer à \textbf{abordagem} própria das \textbf{ciências}, incluindo a investigação, a reflexão, a \textbf{análise} \textbf{crítica}, a imaginação e a criatividade, para investigar causas, elaborar e testar hipóteses, formular e resolver \textbf{problemas} e criar soluções ( inclusive tecnológicas ) com base nos conhecimentos das diferentes áreas. 3.
Repertório cultural - Valorizar e fruir as diversas manifestações artísticas e culturais, das locais às mundiais, e também participar de práticas diversificadas da produção \textbf{artístico-cultural}. 4.
Comunicação - Utilizar diferentes linguagens – verbal ( oral ou \textbf{visual-motora}, como \textbf{Libras}, e escrita ), corporal, visual, sonora e \textbf{digital} –, bem como conhecimentos das linguagens artística, matemática e científica, para se expressar e partilhar informações, experiências, ideias e sentimentos, em diferentes contextos e produzir sentidos que levem ao entendimento mútuo. 5.
Cultura \textbf{digital} - Compreender, utilizar e criar \textbf{tecnologias} \textbf{digitais} de informação e comunicação de forma \textbf{crítica}, significativa, \textbf{reflexiva} e ética nas diversas práticas sociais ( incluindo as escolares ) para se comunicar, acessar e disseminar informações, produzir conhecimentos, resolver \textbf{problemas} e exercer protagonismo e \textbf{autoria} na vida pessoal e coletiva. 6.
Trabalho e projeto de vida - Valorizar a diversidade de saberes e vivências culturais e apropriar -se de conhecimentos e experiências que lhe possibilitem entender as relações próprias do mundo do trabalho e fazer escolhas alinhadas ao exercício da cidadania e ao seu projeto de vida, com liberdade, autonomia, consciência \textbf{crítica} e responsabilidade. 7.
Argumentação - \textbf{Argumentar} com base em fatos, dados e informações confiáveis, para formular, negociar e defender ideias, pontos de \textbf{vista} e decisões comuns, que respeitem e promovam os direitos humanos, a consciência \textbf{socioambiental} e o consumo responsável em âmbito local, regional e global, com \textbf{posicionamento} ético em relação ao cuidado de si mesmo, dos outros e do planeta. 8.
Autoconhecimento e autocuidado - Conhecer -se, apreciar -se e cuidar de sua saúde física e emocional, compreendendo -se na diversidade humana e reconhecendo suas emoções e as dos outros, com autocrítica e \textbf{capacidade} para lidar com elas. 9.
Empatia e cooperação - Exercitar a empatia, o diálogo, a resolução de conflitos e a cooperação, fazendo -se respeitar e promovendo o respeito ao outro e aos direitos humanos, com acolhimento e valorização da diversidade de indivíduos e de grupos sociais, seus saberes, identidades, culturas e potencialidades, sem preconceitos de qualquer natureza. 10.
Responsabilidade e cidadania - Agir pessoal e coletivamente com autonomia, responsabilidade, flexibilidade, resiliência e determinação, tomando decisões com base em princípios éticos, democráticos, inclusivos, sustentáveis e solidários. ( BRASIL, 2018 ) O conceito de \textbf{competência}, adotado pela BNCC, marca a discussão pedagógica e social das últimas décadas e pode ser percebido no texto da LDB ( 1996 ), especialmente quando se estabelecem as finalidades gerais do Ensino Fundamental e do Ensino Médio, nos Artigos 32 e 35.
Além disso, desde as décadas \textbf{finais} do \textbf{século} XX, e ao longo deste início do \textbf{século} XXI, o \textbf{foco} no desenvolvimento de \textbf{competências} tem orientado a maioria dos estados e municípios brasileiros e diferentes países na construção de seus currículos.
É esse também o enfoque adotado nas avaliações internacionais da Organização para a Cooperação e Desenvolvimento Econômico ( OCDE ), que coordena o Programa Internacional de Avaliação de Alunos ( Pisa ), e da Organização das Nações Unidas para a Educação, a \textbf{Ciência} e a Cultura ( Unesco ), que instituiu o Laboratório Latino-americano de Avaliação da Qualidade da Educação para a América Latina ( BRASIL, 2018 ).
Ao adotar essa perspectiva, a BNCC indica que as decisões pedagógicas devem estar orientadas para o desenvolvimento de \textbf{competências}.
Por meio da \textbf{indicação} clara do que os \textbf{alunos} devem “ saber ” ( considerando a constituição de conhecimentos, habilidades, atitudes e valores ) e, sobretudo, do que devem “ saber fazer ” ( considerando a \textbf{mobilização} desses conhecimentos, habilidades, atitudes e valores para resolver \textbf{demandas} \textbf{complexas} da vida cotidiana, do pleno exercício da cidadania e do mundo do trabalho ), a explicitação das \textbf{competências} oferece referências para o fortalecimento de ações que assegurem as aprendizagens essenciais definidas na BNCC ( BRASIL, 2018 ).

% matched lemmas: abordagem, agenda, aluno, análise, argumentar, artístico-cultural, atingir, ator, atual, autoria, capacidade, ciência, competência, complexo, compreensão, crítico, demanda, digital, exigir, final, foco, fomentar, imprescindível, indicação, inter-relacionar, libra, maneira, mobilização, posicionamento, potencial, problema, profundidade, quadro, raciocínio, reflexivo, socioambiental, século, tanto, tarefa, tecnologia, vista, visual-motor
\end{textsample}
